\chapter{C++ pour la latence I~: mémoire}\label{ch:ll:cpp-for-latency-i-memory}

Écrire un octet dans une page de mémoire qui venait d'être mappée a coûté au portable de ce livre environ \qty{1.3}{\micro\second}~:
écrire le même octet dans la même page une seconde fois n'a coûté que \qty{30}{\nano\second}. La première écriture est passée par le noyau,
qui a trouvé une page physique libre, l'a remplie de zéros et l'a installée dans les tables de pages~; la seconde était un simple store ordinaire. Un
processus de trading qui demande de la mémoire au tas pendant qu'il traite un message peut, de temps à autre, emprunter ce chemin, ou plusieurs
autres à l'intérieur de l'allocateur, rapides en moyenne mais lents de loin en loin. Ce chapitre porte sur la mémoire en C++ sur
le chemin critique~: comment les objets sont disposés, ce que coûte une allocation et quand, et les arènes, les réserves et les \omterm{def:ll:cpp-for-latency-i-memory:sbo}{conteneurs à capacité fixe}
qui permettent au chemin critique de fonctionner sans jamais allouer.

\section{Disposition des objets, remplissage et alignement}

\begin{definition}[Remplissage de structure]\label{def:ll:cpp-for-latency-i-memory:padding}
Chaque type scalaire possède un alignement, égal à sa taille sur x86-64 (un \texttt{double} commence à un multiple de 8). Le
compilateur dispose les champs d'une structure dans l'ordre de leur déclaration et insère un
\emph{remplissage de structure}\index{remplissage de structure}, des octets inutilisés, avant tout champ qui serait sinon mal aligné, et à la fin, de sorte que la taille de la structure
soit un multiple de son alignement le plus grand.
\end{definition}

\begin{example}[Ordonner les champs]\label{ex:ll:cpp-for-latency-i-memory:layout}
Un ordre rédigé dans un premier jet déclare un \texttt{char} pour le côté, un \texttt{double} pour le prix, un \texttt{char} pour l'indicateur, un identifiant sur 64 bits,
une quantité sur 32 bits et un \texttt{char} pour la durée de validité~: 23 octets de données occupent 40 octets, dont 17 de remplissage. Triés par taille
(identifiant, prix, quantité, puis les trois caractères), les mêmes champs occupent 24 octets~: huit ordres tiennent dans trois \omterm{def:ll:memory-hierarchy-and-caches:line}{lignes de cache} de 64 octets
au lieu de cinq. Le compilateur ne réordonne pas lui-même les champs~: l'ordre fait partie de l'interface du type.
\end{example}

\begin{omfigure}
\begin{tikzpicture}[font=\footnotesize,x=0.17cm,y=1cm]
\newcommand{\fld}[4]{\draw[fill=#4] (#1,#3) rectangle (#1+#2,#3+0.55);}
\node[anchor=east] at (-1,0.28) {premier jet (40 octets)};
\fld{0}{1}{0}{omDef!45}\fld{1}{7}{0}{gray!15}\fld{8}{8}{0}{omDef!45}\fld{16}{1}{0}{omDef!45}\fld{17}{7}{0}{gray!15}
\fld{24}{8}{0}{omDef!45}\fld{32}{4}{0}{omDef!45}\fld{36}{1}{0}{omDef!45}\fld{37}{3}{0}{gray!15}
\node at (4,0.28) {rempl.}; \node at (12,0.28) {prix}; \node at (20.5,0.28) {rempl.}; \node at (28,0.28) {id}; \node at (34,0.28) {qté};
\node[anchor=east] at (-1,-0.72) {par taille (24 octets)};
\fld{0}{8}{-1}{omMeth!40}\fld{8}{8}{-1}{omMeth!40}\fld{16}{4}{-1}{omMeth!40}\fld{20}{3}{-1}{omMeth!40}\fld{23}{1}{-1}{gray!15}
\node at (4,-0.72) {id}; \node at (12,-0.72) {prix}; \node at (18,-0.72) {qté};
\foreach \x in {0,8,16,24,32,40,48,56,64} {\draw[gray] (\x,-1.25) -- (\x,-1.1); \node[below,font=\scriptsize] at (\x,-1.2) {\x};}
\draw[omThm,dashed] (64,-1.1) -- (64,0.7);
\node[omThm,anchor=south,font=\scriptsize] at (64,0.7) {fin d'une ligne de 64 octets};
\end{tikzpicture}

\omcaption{L'ordre de l'\cref{ex:ll:cpp-for-latency-i-memory:layout} disposé octet par octet~: le premier jet porte 17
octets de remplissage (en gris) sur 40~; les mêmes champs triés par taille en demandent 24, avec le côté, l'indicateur et la durée de validité dans les trois
octets qui suivent la quantité.}\label{fig:ll:cpp-for-latency-i-memory:layout}
\end{omfigure}

Trois règles complémentaires. Placez en tête de structure les champs que lit le chemin critique, de sorte qu'une seule ligne suffise au
cas courant, et déplacez le reste (texte, champs d'audit) dans une structure séparée. Alignez les structures écrites par des threads différents
sur des \omterm{def:ll:memory-hierarchy-and-caches:line}{lignes de cache} distinctes (chapitre~3). Et préférez les indices dans des tableaux aux pointeurs dans les structures chaudes~: un indice sur 32 bits ne pèse que la moitié d'un
pointeur et survit à un déplacement du tableau.

\section{Ce que coûte une allocation, et quand}

\texttt{new} et \texttt{malloc} sont rapides en moyenne. L'allocateur de la bibliothèque C de GNU conserve les petits blocs libérés dans des
caches par thread et les restitue en quelques dizaines de nanosecondes. Il lui arrive aussi, de temps en temps, de faire beaucoup plus~: il regroupe les
blocs libres, agrandit le tas par un appel système, transmet directement au noyau les grands blocs (par défaut à partir de 128~KiB et au-delà, un
seuil qu'il ajuste à mesure que le programme s'exécute), rend la mémoire au noyau lorsque le sommet du tas est libre (par défaut au-delà de
128~KiB), et chaque nouvelle page qu'il touche provoque un \omterm{def:ll:cpp-for-latency-i-memory:fault}{défaut de page}.

\begin{definition}[Défaut de page, pré-chargement des pages]\label{def:ll:cpp-for-latency-i-memory:fault}
Un \emph{défaut de page}\index{défaut de page} est l'interruption par laquelle le processeur passe au noyau lorsqu'un programme touche une page virtuelle qui
n'est pas encore adossée à de la mémoire physique (un défaut mineur~: le noyau alloue et remplit de zéros une page) ou dont le contenu se trouve sur disque
(un défaut majeur). Le \emph{pré-chargement des pages}\index{pré-chargement des pages} consiste à toucher chaque page d'un tampon au démarrage, afin qu'aucun
défaut ne survienne plus tard sur le chemin critique.
\end{definition}

La \Cref{fig:ll:cpp-for-latency-i-memory:faults} mesure les deux coûts. Remplacer l'un de 10\,000 ordres vivants de 64 octets, tiré au hasard,
coûte une paire d'allocation et de libération environ \qty{25}{\nano\second} à la médiane, une réserve un peu moins, et les deux atteignent des dizaines de
microsecondes à leur maximum, lorsque le thread a été interrompu. La première écriture dans chaque page d'un mappage neuf de 64~MiB
coûte plus d'une microseconde par page~: 16\,384 défauts, soit quarante fois le coût des mêmes écritures une fois les pages existantes. Un chemin critique
qui alloue est exposé aux chemins lents de l'allocateur et aux défauts sur une mémoire qu'il n'a pas encore touchée~; celui qui n'alloue pas
n'est exposé ni à l'un ni à l'autre.

\begin{omfigure}
\begin{tikzpicture}
\begin{axis}[name=churn,width=7.2cm,height=5.4cm,ybar,bar width=5pt,ymode=log,log origin y=infty,ymin=5,ymax=200000,
  symbolic x coords={new/delete,pool,pmr pool},xtick=data,enlarge x limits=0.3,title={\small remplacement d'un ordre},
  ylabel={ns (échelle log)},tick label style={font=\footnotesize},label style={font=\small},grid=major,grid style={gray!25},
  legend columns=4,legend style={font=\footnotesize,at={(1.15,-0.22)},anchor=north,draw=none,fill=none,
  /tikz/every even column/.append style={column sep=6pt}},table/col sep=comma]
\addplot[fill=omDef!70,draw=omDef] table[x=allocator,y=p50]{figdata/low-latency/06-cpp-for-latency-i-memory/measured_churn.csv};
\addplot[fill=omDef!35,draw=omDef] table[x=allocator,y=p99]{figdata/low-latency/06-cpp-for-latency-i-memory/measured_churn.csv};
\addplot[fill=omThm!45,draw=omThm] table[x=allocator,y=p999]{figdata/low-latency/06-cpp-for-latency-i-memory/measured_churn.csv};
\addplot[fill=gray!40,draw=gray] table[x=allocator,y=max]{figdata/low-latency/06-cpp-for-latency-i-memory/measured_churn.csv};
\legend{p50,p99,p99.9,maximum}
\end{axis}
\begin{axis}[at={(churn.east)},anchor=west,xshift=2.1cm,width=5.2cm,height=5.4cm,ybar,bar width=14pt,ymode=log,log origin y=infty,
  ymin=5,ymax=5000,symbolic x coords={first touch,pre-faulted},xtick=data,xticklabels={{premier\\ accès},{pré-\\chargée}},xticklabel style={align=center},enlarge x limits=0.5,title={\small écriture d'une page},
  ylabel={ns par page (échelle log)},tick label style={font=\footnotesize},label style={font=\small},grid=major,
  grid style={gray!25},table/col sep=comma]
\addplot[fill=omThm!45,draw=omThm] table[x=kind,y=ns_per_page]{figdata/low-latency/06-cpp-for-latency-i-memory/measured_faults.csv};
\end{axis}
\end{tikzpicture}

\omcaption{À gauche~: destruction et création d'un des 10\,000 ordres vivants de 64 octets, tirés au hasard, avec le tas C++, la réserve du
constructeur et la ressource de réserve de la bibliothèque standard (400\,000 remplacements). À droite~: la première écriture dans chaque page de 4~KiB d'un
mappage neuf de 64~MiB, puis dans les mêmes pages une fois écrites. Mesuré sur un portable (Intel Core Ultra 7 155H) sous WSL2, sans
cœurs isolés, épinglé sur un seul CPU. Données~: \texttt{bench\_alloc.py}.}\label{fig:ll:cpp-for-latency-i-memory:faults}
\end{omfigure}

\section{Arènes, réserves et listes libres}

\begin{definition}[Allocateur par arène]\label{def:ll:cpp-for-latency-i-memory:arena}
Un \emph{allocateur par arène}\index{allocateur par arène} distribue des morceaux consécutifs d'un seul bloc préalloué en faisant avancer
un décalage, et les libère tous d'un coup en le remettant à zéro~; les objets individuels ne sont jamais libérés.
\end{definition}

\begin{definition}[Réserve d'objets, liste libre]\label{def:ll:cpp-for-latency-i-memory:pool}
Une \emph{réserve d'objets}\index{réserve d'objets} préalloue l'espace d'un nombre fixe d'objets d'un même type et le recycle~:
créer un objet prend un emplacement libre, détruire un objet le rend. Une \emph{liste libre}\index{liste libre} est la liste des emplacements libres~;
une liste libre intrusive range ses chaînons à l'intérieur même des emplacements libres, de sorte qu'elle n'a pas besoin de mémoire propre.
\end{definition}

Une arène convient à une mémoire dont la durée de vie est une phase~: l'espace de travail d'un message, reconstruit pour le suivant~; les champs analysés
d'un lot. Une réserve convient à des objets d'un même type aux durées de vie indépendantes~: les ordres, créés et annulés dans n'importe
quel ordre. Les deux ont une capacité fixe décidée à l'avance, et la question de conception qu'elles imposent est la bonne~: que se passe-t-il quand
elle est épuisée. Sur le chemin critique, la réponse n'est jamais «~demander au tas~»~; il faut refuser (rejeter l'ordre, passer en
mode dégradé) et déclencher une alarme, et dimensionner la réserve à partir du pic mesuré, avec une marge.

La bibliothèque standard reprend ces idées sous forme de ressources mémoire (\texttt{std::pmr}, depuis C++17)~: une ressource tampon monotone
est une arène, une ressource de réserve gère des \omterm{def:ll:cpp-for-latency-i-memory:pool}{listes libres} par taille. Elles permettent aux conteneurs standard de les utiliser, au prix d'un appel indirect
à chaque allocation~; la réserve générique de la \cref{fig:ll:cpp-for-latency-i-memory:faults} est plus lente
qu'une réserve pour un seul type.

\section{Petits tampons et conteneurs à capacité fixe}

\begin{definition}[Optimisation des petits tampons, conteneur à capacité fixe]\label{def:ll:cpp-for-latency-i-memory:sbo}
Un type doté d'une \emph{optimisation des petits tampons}\index{optimisation des petits tampons} stocke les petits contenus à l'intérieur de l'objet et
n'alloue que pour les grands~: l'implémentation GNU de \texttt{std::string} conserve jusqu'à 15 caractères en place. Un
\emph{conteneur à capacité fixe}\index{conteneur à capacité fixe} stocke toujours ses éléments à l'intérieur de l'objet, jusqu'à une
capacité fixée à la compilation, et refuse de grandir au-delà.
\end{definition}

L'\omterm{def:ll:cpp-for-latency-i-memory:sbo}{optimisation des petits tampons} fait dépendre l'allocation des données~: un symbole de quatre lettres n'alloue pas, un identifiant
d'ordre client de 22 caractères alloue, et un test avec des identifiants courts ne le verra pas. Le gestionnaire du
\cref{lst:ll:cpp-for-latency-i-memory:naive} est typique des premières versions. Compté avec un
\texttt{operator new}, il alloue neuf fois par message~: une fois pour l'identifiant long, trois fois pendant la croissance du vecteur des exécutions jusqu'à trois, deux fois
pour la nouvelle entrée de la table (le nœud de l'arbre et sa copie de l'identifiant), une fois quand la fonction lambda copie les exécutions, et deux fois quand
\texttt{std::function} stocke la fonction lambda sur le tas et la copie. La réécriture avec une chaîne à taille fixe, un vecteur à taille fixe
des exécutions et une table à taille fixe des entrées fait le même travail sans aucune allocation.

\begin{method}[Supprimer l'allocation d'un chemin critique]\label{met:ll:cpp-for-latency-i-memory:remove}
\begin{enumerate}
\item Compter~: remplacer \texttt{operator new} global dans une compilation de test et vérifier le nombre d'allocations par message après l'échauffement.
\item Pour chaque allocation, décider de la capacité~: identifiants et symboles dans des chaînes à taille fixe, listes par message dans des
vecteurs à taille fixe, objets à durée de vie dans des réserves, espace de travail par message dans une arène réinitialisée à chaque message.
\item Remplacer les objets appelables à effacement de type qui allouent (\texttt{std::function} avec de grandes captures) par des modèles (chapitre~7) ou
par des pointeurs de fonction accompagnés d'un contexte.
\item Allouer et toucher chaque réserve, arène et table au démarrage, dans le fil qui les utilisera (chapitre~4).
\item Garder l'assertion dans les tests, afin qu'une modification ultérieure qui alloue fasse échouer la construction.
\end{enumerate}
\end{method}

\section{Tutoriel~: compter et supprimer les allocations}\label{tut:ll:cpp-for-latency-i-memory:tut}

\begin{tutorial}
\textbf{Objectif.} Compter les allocations d'un gestionnaire de messages, les supprimer et mesurer ce que coûtent l'allocation et les \omterm{def:ll:cpp-for-latency-i-memory:fault}{défauts de page}.
\textbf{État final~:} un gestionnaire sans aucune allocation par message après l'échauffement, et
la \cref{fig:ll:cpp-for-latency-i-memory:faults}.
\begin{tutsteps}
\item \textbf{La première version.} Des chaînes, un vecteur qui grossit, une table, une fonction de rappel.
\omcode{code/low-latency/06-cpp-for-latency-i-memory/cpp/ll_alloc.hpp}{22}{35}{Un gestionnaire de messages qui alloue neuf fois par message.}\label{lst:ll:cpp-for-latency-i-memory:naive}
\item \textbf{Compter.} \texttt{firm\_alloc\_count.hpp} remplace \texttt{operator new} global par une version qui incrémente un
compteur atomique~; le test traite un message, lit le compteur et vérifie 9 pour la première version et 0 pour la réécriture.
\item \textbf{Une réserve avec \omterm{def:ll:cpp-for-latency-i-memory:pool}{liste libre} intrusive.} Chaque emplacement libre contient le pointeur vers l'emplacement libre suivant~; créer
retire un emplacement de la liste, détruire le rajoute.
\omcode{code/firm/arena/cpp/firm_arena.hpp}{45}{80}{La réserve d'objets.}\label{lst:ll:cpp-for-latency-i-memory:pool}
\item \textbf{La même garantie en Rust.} Un allocateur global de comptage (\texttt{firm\_arena::CountingAlloc}) installé
dans le binaire de test vérifie que les opérations de la réserve n'allouent rien.
\item \textbf{Mesurer} avec \texttt{python bench\_alloc.py}~: la rotation des allocations et les \omterm{def:ll:cpp-for-latency-i-memory:fault}{défauts de page}.
\end{tutsteps}
\textbf{Pour aller plus loin.} Réduire l'identifiant client à 15 caractères et recompter~; aligner les emplacements de la réserve
sur les \omterm{def:ll:memory-hierarchy-and-caches:line}{lignes de cache} et comparer les centiles de rotation.
\end{tutorial}

\section{Construction~: la boîte à outils sans allocation}\label{bld:ll:cpp-for-latency-i-memory:arena}

\begin{build}
\bfield{Rôle} Le stockage qu'utilise chaque composant du chemin critique de la partie~IV (la réserve d'ordres du constructeur de carnet, les
états d'ordres de la passerelle, l'espace de travail par message), et le compteur d'allocations sur lequel s'appuie chaque test d'acceptation.
\bfield{Interface} C++20 \texttt{firm::arena}~: \texttt{Arena(bytes)} avec \texttt{allocate(n, align)}, \texttt{make<T>},
\texttt{reset}~; \texttt{Pool<T, N>} avec \texttt{create}, \texttt{destroy}, \texttt{live}~; \texttt{FixedVector<T, N>}~;
\texttt{FixedString<N>}~; \texttt{firm\_alloc\_count.hpp} avec \texttt{allocations()}. Rust \texttt{firm\_arena}~:
\texttt{Arena} (décalages), \texttt{Pool<T>} (indices), \texttt{CountingAlloc} et \texttt{allocations()}.
\bfield{Règles} Les capacités sont fixées à la construction~; l'épuisement renvoie un résultat nul, jamais une allocation sur le tas~; le stockage est
pré-chargé à la construction~; aucune opération après la construction n'alloue.
\bfield{Tests d'acceptation} \texttt{code/firm/arena/}~: décalages alignés, refus lorsque l'arène est pleine et réinitialisation~; ordre de réutilisation de la réserve et
épuisement~; chaînes à taille fixe~; zéro allocation comptée sur l'ensemble, en C++20 et en Rust.
\bfield{Pour aller plus loin} Une réserve dont les emplacements sont alignés sur les \omterm{def:ll:memory-hierarchy-and-caches:line}{lignes de cache} pour des objets écrits par des fils différents~; une arène qui
enregistre son pic d'utilisation pour le dimensionnement.
\end{build}

\begin{omsources}
\item Page de manuel Linux \texttt{mallopt(3)}~: \texttt{M\_MMAP\_THRESHOLD}, \texttt{M\_TRIM\_THRESHOLD}.
\item cppreference, «~\texttt{std::pmr::memory\_resource}~» et ses ressources dérivées.
\item C.~Cook, «~When a microsecond is an eternity~», CppCon 2017 (réutiliser la mémoire sur le chemin critique).
\end{omsources}

\section{Exercices}

\begin{exercise}[$\star$]\label{exo:ll:cpp-for-latency-i-memory:1}
Quelles sont la taille et le remplissage d'une structure contenant, dans cet ordre, un \texttt{char}, un \texttt{int32}, un \texttt{char},
un \texttt{double} et un \texttt{int16}~? Et avec les champs triés par taille~?
\end{exercise}

\begin{exercise}[$\star$]\label{exo:ll:cpp-for-latency-i-memory:2}
Une réserve de 50\,000 ordres de 64 octets~: quelle quantité de mémoire, et combien de pages de 4~KiB faut-il pré-charger~?
\end{exercise}

\begin{exercise}[$\star$]\label{exo:ll:cpp-for-latency-i-memory:3}
Au coût mesuré par page, combien de temps prend le fait de toucher un tampon neuf de 512~MiB au démarrage, et que coûterait-il si
cela se produisait sur le chemin critique, une page à la fois~?
\end{exercise}

\begin{exercise}[$\star\star$]\label{exo:ll:cpp-for-latency-i-memory:4}
Lesquelles de ces instructions allouent avec la bibliothèque standard GNU~: \texttt{std::string s(\textquotedbl{}AAPL\textquotedbl{})}~;
\texttt{std::string s(\textquotedbl{}CLIENT-ORDER-ID-42\textquotedbl{})}~; \texttt{std::vector<int> v; v.reserve(0)}~;
\texttt{std::vector<int> v(1)}~?
\end{exercise}

\begin{exercise}[$\star\star$]\label{exo:ll:cpp-for-latency-i-memory:5}
Le gestionnaire de départ alloue neuf fois par message. Quelles trois modifications suppriment le plus d'allocations, et combien chacune
en supprime-t-elle~?
\end{exercise}

\begin{exercise}[$\star\star$]\label{exo:ll:cpp-for-latency-i-memory:6}
Lire la \cref{fig:ll:cpp-for-latency-i-memory:faults}~: à la médiane, le tas est presque aussi rapide que la réserve. Pourquoi
utiliser malgré tout une réserve~?
\end{exercise}

\begin{exercise}[$\star\star\star$]\label{exo:ll:cpp-for-latency-i-memory:7}
\emph{Programmation.} Ajouter à \texttt{firm::arena::Pool} un contrôle de capacité qui enregistre le high-water mark des objets vivants, lancez avec lui le
benchmark de rotation, et proposez une règle de dimensionnement pour une réserve dont le pic est mesuré sur une année de jours chargés.
\end{exercise}

\begin{exercise}[$\star\star\star$]\label{exo:ll:cpp-for-latency-i-memory:8}
\emph{Trouver la faille.} «~Notre gestionnaire n'alloue rien~: nous l'avons vérifié avec un test qui envoie le même message, le symbole
\texttt{ESZ6} et l'identifiant client \texttt{A1}, un million de fois.~»
\end{exercise}

\section{Problème~: quatre cents microsecondes, une fois par minute}

\begin{problem}[{Problème du week-end --- un gestionnaire qui alloue}]\label{pb:ll:cpp-for-latency-i-memory:1}
Le gestionnaire d'une stratégie prend environ \qty{2}{\micro\second} par message, sauf une ou deux fois par minute où il en prend plusieurs
centaines. L'équipe soupçonne la mémoire. Le gestionnaire est le brouillon du \cref{lst:ll:cpp-for-latency-i-memory:naive}, alimenté par 50\,000
messages par seconde.

\medskip\noindent\textbf{Partie I --- Comptage.}
\begin{enumerate}
\item Combien d'allocations par seconde le brouillon effectue-t-il~?
\item Lister les neuf allocations d'un message et la ligne qui produit chacune.
\item Lesquelles disparaîtraient avec un identifiant de 15 caractères ou moins~?
\item Pourquoi le compte ne dépend-il pas du nombre d'exécutions d'un message au-delà de trois~?
\end{enumerate}

\medskip\noindent\textbf{Partie II --- Les coûts.}
\begin{enumerate}[resume]
\item D'après la \cref{fig:ll:cpp-for-latency-i-memory:faults}, combien coûte une paire d'opérations sur le tas à la médiane~?
\item Combien de temps par seconde les allocations du brouillon coûtent-elles à cette médiane~?
\item Combien coûte un \omterm{def:ll:cpp-for-latency-i-memory:fault}{défaut de page}, et combien en faudrait-il pour porter un message à plusieurs centaines de microsecondes~?
\item Nommer trois événements internes à l'allocateur, rares mais lents.
\end{enumerate}

\medskip\noindent\textbf{Partie III --- La réécriture.}
\begin{enumerate}[resume]
\item Qu'est-ce qui remplace chaque allocation dans le gestionnaire corrigé~?
\item Quelles capacités la réécriture a-t-elle retenues, et que se passe-t-il lorsqu'on en dépasse une~?
\item Pourquoi le stockage du gestionnaire corrigé doit-il être touché au démarrage, et par quel fil~?
\item Combien coûte la recherche linéaire sur 64 entrées de la réécriture, et quand faudrait-il lui substituer une table de hachage~?
\end{enumerate}

\medskip\noindent\textbf{Partie IV --- Le verdict.}
\begin{enumerate}[resume]
\item Énoncer le \emph{résultat nommé}~: allocations par message avant et après la réécriture, et coût d'un \omterm{def:ll:cpp-for-latency-i-memory:fault}{défaut de page}
comparé à une écriture préchargée sur cet ordinateur portable.
\item Comment le test empêche-t-il ce résultat de régresser~?
\item Qu'est-ce qui, outre l'allocation, peut rendre lent un message sur une minute~?
\item Comment montrer quelle cause il s'agissait sur l'hôte de production~?
\item Pourquoi «~refuser quand c'est plein~» est-il acceptable sur le chemin critique, et «~allouer quand c'est plein~» ne l'est-il pas~?
\item Où le jumeau en Rust fait-il respecter la même règle~?
\item Quelle doit être la taille de la réserve si le pic d'ordres vivants sur une année a été de 38\,000~?
\item En une phrase~: quelle est la règle mémoire d'un chemin critique~?
\end{enumerate}
\end{problem}

\section{Questions d'entretien}

\begin{interviewq}[$\star$ \iqroles{développeur}]\label{iq:ll:cpp-for-latency-i-memory:1}
Pourquoi réordonner les champs d'une structure peut-il rendre un programme plus rapide~?
\end{interviewq}

\begin{interviewq}[$\star\star$ \iqroles{développeur}]\label{iq:ll:cpp-for-latency-i-memory:2}
Pourquoi les systèmes à faible latence évitent-ils l'allocation sur le tas dans le chemin critique, si \texttt{malloc} ne prend que 20 nanosecondes~?
\end{interviewq}

\begin{interviewq}[$\star\star$ \iqroles{développeur}]\label{iq:ll:cpp-for-latency-i-memory:3}
Implémentez une \omterm{def:ll:cpp-for-latency-i-memory:pool}{réserve d'objets} pour un type de taille fixe. Que se passe-t-il lorsqu'elle est vide~?
\end{interviewq}

\begin{interviewq}[$\star\star$ \iqroles{développeur}]\label{iq:ll:cpp-for-latency-i-memory:4}
Qu'est-ce que l'optimisation des petites chaînes, et pourquoi peut-elle masquer des allocations à un test~?
\end{interviewq}

\begin{interviewq}[$\star\star$ \iqroles{développeur}]\label{iq:ll:cpp-for-latency-i-memory:5}
Qu'est-ce qu'un \omterm{def:ll:cpp-for-latency-i-memory:fault}{défaut de page}, et comment les tenez-vous à l'écart d'un chemin critique~?
\end{interviewq}

\begin{interviewq}[$\star\star\star$ \iqroles{développeur}]\label{iq:ll:cpp-for-latency-i-memory:6}
Comment prouveriez-vous, en intégration continue, qu'un composant n'alloue jamais après le démarrage~?
\end{interviewq}
